\section{Certified removal of visible boundary layers}
\label{sec:boundary-unlayer}

The collar construction in Section~\ref{sec:boundary-transversal}
suggests a concrete preprocessing operation: remove an exposed outer
layer before building the Euler search model.  The maintained optional
\path{unlayer_boundary=True} performs this operation with an independent
certificate.  It preserves the compact manifold and decreases the
tetrahedron count.  It neither preserves the labelled normal solution
cone nor supplies a bound on its resolving number $h$.

The operation is restricted to finite, connected, orientable,
one-vertex sources with one torus boundary.  These are the sources
produced by the vertex reduction of Section~\ref{sec:vertex-reduction}
unless an unknot proof has already terminated the chain.  Sources with
more vertices are left unchanged by this optional operation.  All
normal search methods can use the resulting smaller source; the default
\path{unlayer_boundary=False} retains the previous preprocessing policy.

\subsection{The unique visible candidate and its labelled deletion}

A one-vertex triangulation of the torus has two boundary triangles.
If they belong to different tetrahedra, there is no exposed layer of
the supported kind.  Otherwise their common tetrahedron $d$ is the only
candidate.  It has two boundary faces and two attached faces.  Both
attached faces must lead to surviving tetrahedra, and their target
face records must be distinct.  The two targets may lie in the same
surviving tetrahedron.  Self-pairings involving $d$ and deletion of the
only tetrahedron are excluded.

Let $T$ be the supplied source and $S$ the proposed successor.  Delete
$d$, expose its two reciprocal attachment faces, and keep every other
face pairing.  Surviving tetrahedra retain their order in the original
label list and are compacted to consecutive labels; all local vertex
labels remain unchanged.  Thus $S$ is uniquely determined by $T$ and
$d$.  The checker compares every supplied successor face record with
this literal deletion.  It separately validates both sources as compact
orientable manifolds with one vertex and one torus boundary.  The two
boundary faces of $S$ must be exactly the exposed attachment targets.

These conditions allow identifications among the tetrahedron's vertices
and edges.  In particular, requiring four distinct global vertices would
exclude every intended one-vertex example.  The earlier embedded
boundary-shelling checker has that requirement and remains unchanged.
The present operation uses the generalized boundary-layering theorem
instead.

\subsection{Reconstructing the forward layering independently}

Write $f_0,f_1$ for the attached face indices in $d$, and put
\[
 \{a,b\}=\{0,1,2,3\}\setminus\{f_0,f_1\},\qquad a<b.
\]
The attached faces share the local edge $ab$.  Map this ordered edge
through each attachment permutation into $S$.  The two occurrences
must represent the same global boundary edge $e$, with the same
oriented ends.  Equality of the unique global vertex alone would not
check this condition: a loop edge still has two distinguished oriented
ends.  The checker reconstructs edge roots and orientation parities
from the original face pairings of $S$ to compare the two maps.

Next map every surviving local edge incidence of $S$ back into $T$.
All occurrences of a global edge must give the same target edge and
the same relative orientation.  Distinct edge classes of $S$ must have
distinct images.  Hence adding $d$ does not merge existing edge classes
or reverse their end identifications.  There must be exactly one extra
edge in $T$, namely the opposite local edge $f_0f_1$ of $d$, denoted
$e'$.  If $\iota$ is the verified map on old edge classes, the boundary
edge sets must satisfy the exact equality
\[
 E_\partial(T)
   =\bigl(\iota(E_\partial(S))\setminus\{\iota(e)\}\bigr)
       \cup\{e'\}.
\]
The total edge count and the set of all surviving edge images are also
checked.  The four side edges of the attaching quadrilateral need not
be distinct; the comparisons use their reconstructed global classes.

\begin{proposition}[Soundness of a certified inverse layer]
\label{prop:boundary-unlayer-sound}
An accepted inverse-layer packet identifies a labelled triangulation
$S$ with one fewer tetrahedron than $T$, representing the same compact
manifold.  Its successor is nonempty.  In particular, if $T$ is an
authenticated classical knot exterior, then $S$ is an authenticated
exterior of the same knot.
\end{proposition}
\begin{proof}
The literal deletion comparison reconstructs $T$ by adding the fresh
tetrahedron $d$ to $S$ along two distinct boundary triangles.  The
shared-edge comparison proves that both face permutations extend the
same oriented identification of their joining edge.  This is the
layering operation of Jaco--Rubinstein, \emph{Layered-triangulations of
3-manifolds}, Section~3, printed page~9; see also Lemma~9.2, printed
pages~85--86, in
\url{https://arxiv.org/pdf/math/0603601v2}.
Their statement allows identifications among the quadrilateral's sides
and produces a triangulation of the same compact manifold.  It applies
without assuming an embedded attaching disc.  The independent manifold
and edge-incidence checks exclude the unsupported degeneracies above.
The source reconstruction removes exactly one tetrahedron and explicitly
requires a nonempty successor, proving the remaining assertions.
\end{proof}

The packet records an unmarked compact-manifold move.  Recognition here
uses the authenticated knot-exterior type.  A future consumer requiring
a specified boundary slope or cocycle would also have to transport that
data through the checked diagonal flip; unchanged local labels alone
are not a peripheral-coordinate map.

\subsection{A small certificate and polynomial preprocessing}

The public source-only constructor is
\path{boundary_unlayer.unlayer_boundary_once}.  It returns either
\path{UNLAYERED} with a packet, or \path{UNCHANGED}.  The independent
entry point is \path{boundary_unlayer_verify.verify_boundary_unlayering}.
It imports the source-incidence validator and no producer, search
backend, LP solver, external topology engine, or research module.
The packet has exactly four fields:
\begin{center}
\begin{tabular}{ll}
\path{schema} & \path{normal-boundary-unlayer-v1} \\
\path{source_sha256} & canonical JSON digest of the complete labelled $T$ \\
\path{tetrahedron} & literal integer label of $d$ in $T$ \\
\path{successor} & the complete densely labelled $S$
\end{tabular}
\end{center}
The source digest uses sorted-key JSON with fixed compact separators.
The topology check still uses the entire supplied source.  Boolean
tetrahedron labels and nonintegral pairing entries are rejected.
Malformed packets return false; caller exceptions propagate with their
original identity, including exceptions raised during either source
validation.

\begin{proposition}[Polynomial visible-collar removal]
\label{prop:boundary-unlayer-polynomial}
Starting with $t$ tetrahedra, repeated accepted inverse layers perform
at most $t-1$ removals.  Candidate discovery, certificate construction
and independent replay use polynomial combinatorial bit work and
polynomial certificate length.
\end{proposition}
\begin{proof}
Each accepted step strictly decreases the positive tetrahedron count.
There is at most one candidate at each stage, determined by the two
boundary face incidences, so no exponential family of candidate
subsets is inspected.

At size $s$, the face-record encoding has $O(s\log(s+2))$ bits: each
face has one tetrahedron label and a permutation of four vertices.
The shared source validator processes $O(s)$ face, edge and vertex
incidences.  Union by size bounds any union--find parent path by
$O(\log(s+2))$.  All indices and incidence counters have
$O(\log(s+2))$ bits; the sparse matching rows it constructs have
constant-size coefficients.  Connectivity, link and boundary checks
process only polynomially many such data.  The unlayer checker adds
literal face comparisons, one source digest and maps on $O(s)$ edge
incidences.  Even allowing linear scans for each associative lookup,
\[
 P(s)=O\bigl(s^2\log^2(s+2)\bigr)
\]
is a conservative bit bound for a source check and a complete candidate
attempt.  The sum of these bounds over at most $t$ stages is
$O(t^3\log^2(t+2))$.  This deliberately loose bound needs no average
hash-table assumption.  Storing every labelled successor explicitly
uses $O(t^2\log(t+2))$ bits.  Repeating the independent checks during
chain replay changes these polynomial bounds only by a constant factor.
\end{proof}

Successive removals can reverse a collar built by explicit forward
layerings, working from its exposed outermost tetrahedron inward.
This does not establish a unique simplified triangulation or recognize
every possible collar presentation.  Further legal removals from the
earlier source are also allowed; every accepted move has its own
homeomorphism justification.

\subsection{Versioned chains, shared limits and fresh source algebra}

The original \path{diagram-crushing-v1} chain retains its strict
normal-crush-only interpretation.  The additive
\path{diagram-crushing-v2} chain permits both unchanged
\path{normal-ray-crushing-v1} packets and the new inverse-layer packets.
Verification authenticates the diagram exterior and optional original
shelling prefix, then dispatches each step to its independent checker.
Every nonterminal step strictly decreases the tetrahedron count.
An inverse layer always has a nonempty successor; only a verified
terminal normal crush can end an unknot chain.

When enabled, the driver tries the unique inverse-layer candidate after
vertex reduction and before building the next Euler search model.
An accepted step restarts the source loop.  All Q models, coordinates,
dual caches, learned conflicts and final-source exclusion packets belong
to the preceding labelled source and are rebuilt when needed.  No claim
that a normal ray or its conflict set survives the move is used.
Knottedness still requires complete independently checked Euler coverage
on the final source, after authentication of the mixed chain.  A partial
chain remains a \path{REDUCED} source certificate, not a recognition
verdict.

With this option enabled, \path{max_crushes} charges both kinds of
strict reduction.  Its unit is the number of mixed-chain steps; the
statistics \path{crushing_steps} and \path{boundary_unlayerings}
separately count normal crushes and inverse layers.  The default false
option preserves the previous unit and result schema.  All LP pivots,
anchors, candidate work, tree nodes, ellipsoid updates and feasibility
calls already spent remain charged across these source changes.
Limit exhaustion produces \path{INCONCLUSIVE} with its partial chain.
Candidate validation itself is polynomial preprocessing outside the
search allowances.

Visible-collar removal directly addresses the construction that moves
boundary incidences away from a fixed core in the preceding section.
The normal cone after removal may change, and neither $h$ nor $\delta$
is asserted to be invariant or monotone.  Sources with their two
boundary faces in different tetrahedra may still have difficult conflict
hypergraphs.  A uniform small-parameter theorem remains a separate
requirement for general quasi-polynomial recognition.
